
%===================================
%===================================
\section{Conclusiones}
\label{sec:Conclusiones}


\begin{enumerate}

    \item La implementación secuencial optimizada por compilador alcanza entre
    $11.6\times$ y $14.5\times$ de speedup sobre la implementación secuencial
    base, mientras que ninguna configuración paralela supera $0.74\times$. Esto
    se debe a que las optimizaciones de compilador activan la vectorización AVX2,
    que permite procesar cuatro valores en paralelo dentro del mismo núcleo, sin
    ningún costo de sincronización. Antes de recurrir a la paralelización, vale
    la pena asegurarse de que el código serial esté bien optimizado.

    \item El algoritmo de Jacobi 1D no se beneficia de la paralelización para
    los tamaños de problema evaluados: ninguna implementación paralela supera a
    la versión secuencial base. La razón es que el cálculo del residuo se
    ejecuta en un solo hilo mientras los demás esperan, y ese trabajo tiene el
    mismo costo que el que realiza cada hilo en su porción del arreglo,
    $\mathcal{O}(n/p)$. Según la Ley de Amdahl, esta fracción serial limita el
    speedup máximo alcanzable sin importar cuántos hilos se agreguen.

    \item A medida que $n$ crece, el cómputo útil por iteración empieza a
    compensar el costo fijo de sincronización. Para la implementación con hilos
    y $p=6$, el speedup pasa de $0.073\times$ en $n=100$ a $0.726\times$ en
    $n=2000$. Aun así, con los tamaños evaluados nunca se supera la línea
    de~$1\times$.

    \item La configuración óptima para la implementación con hilos es $p=6$,
    lo que coincide exactamente con los seis núcleos físicos del procesador
    utilizado. Al pasar a $p=8$, el speedup cae: para $n=1000$,
    baja de $0.487\times$ a $0.304\times$. Los hilos adicionales comparten el
    mismo núcleo físico y compiten por los mismos recursos de ejecución, por lo
    que el Hyperthreading no aporta ningún beneficio en este caso.

    \item Para $p=2$, hilos y procesos producen tiempos muy similares en todos
    los valores de $n$. A partir de $p=4$, los procesos empiezan a ser
    considerablemente más lentos, y la brecha se amplía con $p$. El caso más
    notable se da en $n=1000$ con $p=12$: la implementación con hilos tarda
    $56{,}775$~ms frente a $118{,}164$~ms de la implementación con procesos,
    es decir, $2.1\times$ más lento.

    \item La implementación secuencial con alineación explícita de memoria a
    64~bytes no muestra ninguna mejora significativa respecto a la versión base,
    con speedups entre $0.988\times$ y $1.039\times$. Para los tamaños de
    problema evaluados, los arreglos caben en los niveles L2 y L3 de caché, y
    el patrón de acceso secuencial es suficiente para que el procesador gestione
    los datos de forma eficiente por sí solo.

\end{enumerate}